\documentclass[11pt]{article}

\usepackage{amssymb,amsmath,amsfonts,eurosym,geometry,ulem,graphicx,caption,color,setspace,sectsty,comment,footmisc,natbib,pdflscape,subfigure,array,hyperref}
\usepackage{dsfont,amsthm}
\hypersetup{hidelinks}

\normalem

\singlespacing
\newtheorem{theorem}{Theorem}
\newtheorem{corollary}[theorem]{Corollary}
\newtheorem{proposition}{Proposition}
\newtheorem{lemma}{Lemma}
\newtheorem{claim}{Claim}

\newtheorem{hyp}{Hypothesis}
\newtheorem{subhyp}{Hypothesis}[hyp]
\renewcommand{\thesubhyp}{\thehyp\alph{subhyp}}

\newcommand{\red}[1]{{\color{red} #1}}
\newcommand{\blue}[1]{{\color{blue} #1}}

\newcolumntype{L}[1]{>{\raggedright\let\newline\\\arraybackslash\hspace{0pt}}m{#1}}
\newcolumntype{C}[1]{>{\centering\let\newline\\\arraybackslash\hspace{0pt}}m{#1}}
\newcolumntype{R}[1]{>{\raggedleft\let\newline\\\arraybackslash\hspace{0pt}}m{#1}}

\geometry{left=1.0in,right=1.0in,top=1.0in,bottom=1.0in}
\setlength{\parskip}{0.7em plus 0.15em minus 0.1em}
\setlength{\parindent}{0pt}

\renewenvironment{abstract}
  {\small\begin{center}\bfseries\abstractname\end{center}\noindent}
  {\par}

\begin{document}

\title{Pseudo-Rational Choice with Multiple Tastes and Beliefs}
\author{Gil Riella\thanks{Department of Economics, University of Brasília} \and Leonardo Sampaio\thanks{Department of Economics, University of Brasília}}
\date{\today}
\maketitle
\begin{abstract}
This paper aims to provide an axiomatic characterization of pseudo-rational
multi-taste, multi-prior choice over menus of Anscombe--Aumann acts. Taking a
choice correspondence as primitive, the framework allows for violations of
the Weak Axiom of Revealed Preference. The main objective is to characterize
choice correspondences admitting a representation of the form
\[
    c(A)=\bigcup_{(\pi,u)\in\mathcal{M}}
    \operatorname*{arg\,max}_{f\in A}
    \sum_{s\in S}\pi(s)\mathbb{E}_{f(s)}[u],
\]
where $\mathcal{M}$ is a set of prior--utility pairs, with each prior
representing beliefs over states and each utility function representing a
possible evaluation of prizes.

The model follows the notion of justifiable choice: an act is chosen whenever
it maximizes subjective expected utility under at least one admissible belief--taste
pair. This allows for uncertainty regarding both the states of the world and
the decision maker's tastes.

The analysis focuses on compact metric prize spaces. The proposed approach draws on
maximal menu arguments and Heller's (2012) minimax theorem method to
investigate the connection between revealed choice and multi-utility
representations. Finally, the paper seeks to explore how indecisiveness and
indifference may be identified and distinguished from observed choice
behavior.

\medskip
\noindent\textbf{Keywords:} Menu effects, Multiple priors, Multiple utilities, Incomplete preferences, Revealed preference, Indecisiveness.

\medskip
\noindent\textbf{JEL Codes:} D81, D83
\end{abstract}

\newpage

\section{Introduction} \label{sec:introduction}

The literature on incomplete preferences provides a natural framework for studying decision makers who are unable to rank every pair of alternatives. In environments under uncertainty, such incompleteness may arise from indecisiveness about tastes, beliefs, or both. A substantial literature studies these possibilities by taking a binary relation as primitive. In particular, Nau (2006), Galaabaatar and Karni (2013), Riella (2015), and Hara and Riella (2023) provide different axiomatic foundations for representations in which preferences are described by multiple probability--utility pairs. Hara and Riella (2023) extend this literature to an arbitrary compact metric prize space, allowing both tastes and beliefs to be simultaneously imprecise.

A different question arises when the primitive object is not a possibly incomplete preference relation, but a choice correspondence. In this case, indecisiveness must be inferred directly from observed choices across menus. Eliaz and Ok (2006) provide an important starting point by showing how choice behavior can distinguish indifference from indecisiveness and by introducing WARNI as a central consistency condition. Subsequently, Heller (2012) develops a justifiable choice model that allows for indecisiveness either in tastes or in beliefs, depending on the environment considered, but not simultaneously in both dimensions. Implicitly, he does so through a pseudo-rational representation in the sense of Aizerman and Malishevski (1981), where the decision maker behaves as if endowed with a collection of preferences, and an alternative is chosen whenever it is maximal according to at least one of them.

Following a related idea, Safonov (2023) studies ambiguity through the lens of framing while also taking a choice correspondence as primitive. In his model, the analyst observes the set of choices that may arise under different frames. Each frame determines a set of priors, while the utility function is common across frames, and conditional on a given frame the decision maker behaves as a maxmin expected utility maximizer. Hence, observed choice can be understood as the union of the choices generated by the different selves available to the decision maker.

Our main goal is to obtain a pseudo-rational representation in which the individual behaves as if endowed with a collection of Anscombe--Aumann preferences. Equivalently, we focus on obtaining a representation analogous to that of Hara and Riella (2023), but taking a choice correspondence as primitive while following a similar general strategy. First, in Theorem 1, we obtain a representation in terms of state-dependent utilities and then impose a variant of a strict monotonicity condition in order to recover the appropriate state-independent structure from extreme or exposed elements of the auxiliary convex representation. As emphasized by Nau (2006), and subsequently developed by Riella (2015) and Hara and Riella (2023), this geometric step is what allows a state-dependent representation to be decomposed into probabilities over states and state-independent utility functions.

Our intended contribution is both behavioral and technical. From the behavioral perspective, we aim to provide a clearer characterization of how simultaneous indecisiveness in tastes and beliefs can be identified from observed choice behavior. Once this general representation is obtained, we also intend to characterize the conditions under which indecisiveness can be restricted to only one of these dimensions. More specifically, we are interested in identifying which additional assumptions on the choice correspondence lead to a single-prior expected multi-utility representation, and which lead to a single-taste multi-prior representation. In this sense, the general multi-prior, multi-taste model provides a common framework from which these two important special cases can be derived.

From the technical perspective, we work with a compact metric prize space rather than restricting attention to a finite set of prizes. This follows the extension developed by Hara and Riella (2023) and allows the model to cover substantially richer environments. In addition, we avoid imposing the existence of best and worst prizes, an assumption that was common in earlier axiomatizations of multiple tastes and beliefs. Hara and Riella (2023) already show that this restriction can be weakened in the preference-based environment, and our objective is to obtain an analogous flexibility while working directly with a choice correspondence.

\section{Setup} \label{sec:result}

We consider an Anscombe--Aumann setup with a finite state space $S$ and a compact metric prize space $X$. An act is a function that maps each state of the world $s\in S$ to a lottery $f(s)\in\Delta(X)$, and we denote the set of all acts by $\mathcal{F}$.

Here, $\Delta(X)$ denotes the set of Borel probability measures on $X$, endowed with any metric that metrizes the topology of weak convergence, and $\mathcal{F}:=\Delta(X)^S$ is endowed with the product metric. Thus, $f_n\to f$ in $\mathcal{F}$ if and only if $f_n(s)$ converges weakly to $f(s)$ for every $s\in S$. Since $X$ is compact and metrizable and $S$ is finite, both $\Delta(X)$ and $\mathcal{F}$ are compact metrizable spaces. Whenever continuous utility functions are considered, $C(Y)$ denotes the space of continuous real-valued functions on a compact metric space $Y$, endowed with the supremum norm. In particular, state-dependent utility functions belong to $C(X\times S)$, while state-independent utility functions belong to $C(X)$.

Let $\mathcal{A}$ be the collection of nonempty compact menus of acts, endowed with the Hausdorff metric induced by the metric on $\mathcal{F}$, so that $\mathcal{A}\subseteq 2^{\mathcal{F}}\setminus\{\emptyset\}$. Our primitive is a choice correspondence $c\colon\mathcal{A}\to 2^{\mathcal{F}}\setminus\{\emptyset\}$ such that, for every $A\in\mathcal{A}$, $\emptyset\neq c(A)\subseteq A$. The choice correspondence is an object that the analyst can usually observe and is therefore more primitive than the notion of preference (Ok and Eliaz, 2006). Its interpretation is as follows: when the decision maker is presented with a menu $A$, she may choose any act in $c(A)$.

We define mixtures of acts state by state. That is, for every $\alpha\in[0,1]$, every $g,h\in\mathcal{F}$, and every $s\in S$, $\bigl(\alpha g+(1-\alpha)h\bigr)(s)=\alpha g(s)+(1-\alpha)h(s)$. Hence, $\mathcal{F}$ can also be viewed as a mixture space. We abuse notation by using $x\in X$ to denote the degenerate constant act that yields outcome $x$ with certainty in every state, and $p\in\Delta(X)$ to denote the constant act that yields lottery $p$ in every state $s\in S$. For $\nu\in C(X\times S)$ and $f\in\mathcal{F}$, the associated state-dependent expected utility is
\[
    U_\nu(f):=\sum_{s\in S}\int_X \nu(x,s)\,df(s)(x).
\]
For a prior $\pi\in\Delta(S)$ and a state-independent utility function $u\in C(X)$, define the subjective expected utility of $f\in\mathcal{F}$ by
\[
    V_{\pi,u}(f):=\sum_{s\in S}\pi(s)\int_X u(x)\,df(s)(x).
\]
Here, $\Delta(S)$ is the simplex of probability distributions on $S$, endowed with its Euclidean topology. The space $\Delta(S)\times C(X)$ carries the product of this topology and the supremum-norm topology. Since $u$ is continuous and $X$ is compact, the integrals are well defined and $V_{\pi,u}$ is continuous and affine on $\mathcal{F}$. For a constant act $p\in\Delta(X)$, $V_{\pi,u}(p)=\int_X u(x)\,dp(x)$.

\section{Representation Theorems}

In this section, we shall present our main results, which consist of the
axiomatization of the additively separable expected multi-utility and the
multi-prior expected multi-utility representations when the set of prizes
is a compact metric space. We also discuss the axioms used, the intuition
for the proofs (which arise from interesting facts implied by the axioms),
and what the theorems mean behaviorally.

\subsection{The Additively Separable Expected Multi-Utility Representation}

We impose the following axioms on the choice correspondence $c$:

\medskip
\noindent\textbf{Axiom 1 (Pseudo-WARP).} For all $A,B\in\mathcal{A}$, if $f\in c(B)\cap A$ and $c(A)\subseteq B$, then $f\in c(A)$.

\medskip
\noindent\textbf{Axiom 2 (Independence).} For every $A\in\mathcal{A}$, $\lambda\in(0,1)$, and $h\in\mathcal{F}$,\footnote{We abuse notation and write $\lambda A+(1-\lambda)h$ for the set $\{\lambda f+(1-\lambda)h:f\in A\}$.} $c\bigl(\lambda A+(1-\lambda)h\bigr)=\lambda c(A)+(1-\lambda)h$.

\medskip
\noindent\textbf{Axiom 3 (Continuity).} The set $\bigl\{(A,f)\in\mathcal{A}\times\mathcal{F}:f\in c(A)\bigr\}$ is closed.

\medskip
\noindent\textbf{Axiom 4 (Indirect Ambiguity Aversion).} For every $A\in\mathcal{A}$, $\lambda\in[0,1]$, $f\in A$, and $g,h\in\mathcal{F}$, if $f\notin c\bigl(A\cup\{h\}\bigr)$ and $f\notin c\bigl(A\cup\{g\}\bigr)$, then $f\notin c\bigl(A\cup\{\lambda g+(1-\lambda)h\}\bigr)$.

\medskip
\noindent\textbf{Axiom 5 (No-Hedging).} For every $A\in\mathcal{A}$, $\lambda\in[0,1]$, and $g,h\in\mathcal{F}$, if $\{g,h\}\subseteq A$, then $c(A)\subseteq c\bigl(A\cup\{\lambda g+(1-\lambda)h\}\bigr)$.

\medskip
\noindent\textbf{Axiom 6 (Constant Non-Triviality).} There exist
$\hat p,\hat q\in\Delta(X)$ such that
$\hat q\notin c\bigl(\{\hat p,\hat q\}\bigr)$.

\medskip
Pseudo-WARP imposes consistency across menus while allowing the choice correspondence to be generated by multiple possible rationales. Independence requires choice to commute with a common mixture transformation, and Continuity rules out discontinuous changes in choice as menus and acts converge. Indirect Ambiguity Aversion imposes a form of consistency under hedging: if $f$ is not chosen in the presence of $h$, and $h$ is not chosen in the presence of $g$, then replacing $h$ with any mixture of $g$ and $h$ cannot make $f$ chosen. 

No-Hedging requires that adding a mixture of two acts already available in a menu does not eliminate any act that was previously chosen. Constant Non-Triviality requires that there exist two constant acts $\hat p,\hat q\in\Delta(X)$ such that $\hat q$ is not chosen from $\{\hat p,\hat q\}$.

Our choice-based analysis requires an auxiliary representation of a strict
partial order by a compact, convex set of continuous state-dependent
utilities. While Galaabaatar and Karni (2012) study related representations in
finite-dimensional spaces, our setting allows a compact metric prize space.
Proposition~\ref{prop:strict-representation} establishes the required result by adapting the approach of
Hara and Riella (2023). Consider a strict relation $\succ\,\subseteq\mathcal{F}\times\mathcal{F}$
and the following axioms:

\medskip
\noindent\textbf{(S1) Independence.} For all $f,g,h\in\mathcal{F}$ and
$\alpha\in(0,1]$, $f\succ g
    \quad\Longleftrightarrow\quad
    \alpha f+(1-\alpha)h\succ\alpha g+(1-\alpha)h$.

\medskip
\noindent\textbf{(S2) Open Continuity.} For every $f\in\mathcal{F}$,
the sets $\{g\in\mathcal{F}:f\succ g\}$ and
$\{g\in\mathcal{F}:g\succ f\}$ are open in $\mathcal{F}$.

\medskip
\noindent\textbf{(S3) Constant Non-Triviality.} There exist
$\hat p,\hat q\in\Delta(X)$ such that $\hat p\succ\hat q$, where lotteries
are identified with the corresponding constant acts.

\medskip
\begin{proposition}\label{prop:strict-representation}
A strict partial
order $\succ\,\subseteq\mathcal{F}\times\mathcal{F}$ satisfies
(S1)--(S3) if and only if there exists a nonempty compact and convex set
$\mathcal{V}\subset C(X\times S)$ and there exist
$\hat p,\hat q\in\Delta(X)$ such that $U_\nu(\hat p)>U_\nu(\hat q)
 \qquad\text{for every }\nu\in\mathcal{V}$,
and, for every $f,g\in\mathcal{F}$,
\[
    f\succ g
    \quad\Longleftrightarrow\quad
    U_\nu(f)>U_\nu(g)
    \quad\text{for every }\nu\in\mathcal{V}.
\]
\end{proposition}



The proof, given in Appendix~\ref{sec:appendix}, follows the same path as in Theorem 8 of Hara and Riella (2023), but removes the axiom of One-Side Monotonicity over the strict partial order. That allows the definition of an appropriate preorder to apply Theorem 1 from the previously mentioned paper. Building on Proposition~\ref{prop:strict-representation}, we can follow Heller's (2012) approach and apply a minimax theorem to obtain a justifiable choice representation. The argument begins by constructing an auxiliary relation from the maximal menus associated with each act in the Anscombe--Aumann framework.

\begin{theorem}\label{thm:justifiable-choice}
The following statements are equivalent:
\begin{enumerate}
    \item[(1)] The choice correspondence $c$ satisfies Axioms 1--6.
    \item[(2)] There exist distinct lotteries $\hat p,\hat q\in\Delta(X)$ and
    a nonempty compact set $\mathcal{W}\subseteq C(X\times S)$ such that,
    for every $A\in\mathcal{A}$,
    \[
        c(A)=\bigcup_{\nu\in\mathcal{W}}
        \operatorname*{arg\,max}_{g\in A} U_\nu(g),
    \]
    and, for every $\nu\in\mathcal{W}$,
    $U_\nu(\hat p)>U_\nu(\hat q)$.
\end{enumerate}
\end{theorem}

\textit{Intuition for the proof.} We provide the intuition for the sufficiency direction. First, we begin with binary choice. Define $f\succ_c g$ whenever $g\notin c(\{f,g\})$. Under our axioms, this is a strict partial order satisfying the conditions of Proposition 1, and therefore it can be represented by a compact and convex family $\mathcal V$ of continuous state-dependent utilities. Thus, whenever $g$ is not revealed strictly better than $f$, there is some $\nu\in\mathcal V$ under which $f$ performs at least as well as $g$.

Next, fix an arbitrary choice $f\in c(A)$. We use Zorn's Lemma, as in Safonov (2023), obtaining a maximal menu $M\supseteq A$ such that $f\in c(M)$. Intuitively, $M$ contains all alternatives that can coexist with $f$ without dislodging it from choice. In addition to this, we mix $M$ and $f$ with a constant full-support act $q$, obtaining $f_\lambda\in c(M_\lambda)$, and then extend $M_\lambda$ to another maximal menu $M^*$.

Additionally, our axioms imply that both $M^*$ and $\mathcal F\setminus M^*$ are convex. This geometry is crucial. From binary choice, for every $g\in M^*$ there exists some $\nu\in\mathcal V$ such that $U_\nu(f_\lambda)\geq U_\nu(g)$. The difficulty is that this utility may depend on $g$. The minimax theorem then allows us to reverse these quantifiers and obtain a single $\nu^*\in\mathcal V$ such that $U_{\nu^*}(f_\lambda)\geq U_{\nu^*}(g)$ for every $g\in M^*$. Thus, one utility supports $f_\lambda$ against the entire maximal menu.

Moreover, the convexity of $M^*$ and of its complement allows us to strengthen this supporting property. We show that $M^*$ is exactly the lower contour set generated by $f_\lambda$,
\[
 M^*=\{g\in\mathcal F:U_{\nu^*}(g)\leq U_{\nu^*}(f_\lambda)\},
\]
and, correspondingly, $M=\{g\in\mathcal F:U_{\nu^*}(g)\leq U_{\nu^*}(f)\}$. Hence, geometrically, the boundary of a maximal menu is a level set of $U_{\nu^*}$. The full-support transformation additionally allows us to find acts lying strictly on each side of this boundary and to show that every act on the boundary is chosen.

Finally, take any menu $B$ and any $U_{\nu^*}$-maximizer $g$ in $B$. By mixing the entire menu with an appropriate act, we move $g$ exactly onto the boundary of $M^*$, while all other alternatives remain weakly inside it. Since boundary acts are chosen, Pseudo-WARP and Independence allow us to transport this choice back to the original menu, giving $g\in c(B)$.

We conclude by defining $\mathcal W\subseteq\mathcal V$ as the family of utilities with this property. The previous argument shows that every observed choice is justified by some element of $\mathcal W$, while every maximizer of an element of $\mathcal W$ is chosen. Finally, closedness of $\mathcal W$ in the compact set $\mathcal V$ gives compactness and completes the desired representation.

\subsection{Multi-Prior Expected Multi-Utility Representation}

We now axiomatize the multi-prior expected multi-utility representation
using a choice correspondence as the primitive. Galaabaatar and Karni
(2013) axiomatize this model using a strict partial order, with related
approaches developed by Riella (2015) and Hara and Riella (2023).
Choice-based approaches, notably Heller (2012) and Safonov (2023), do not
simultaneously allow for multiple priors and multiple utility functions.

\medskip
\noindent\textbf{Axiom 7 (Strict One-Side Monotonicity).} For every $A\in\mathcal{A}$ and every $f,g\in\mathcal{F}$, if $f\notin c\bigl(A\cup\{g(s)\}\bigr)$ for every $s\in S$, then $f\notin c\bigl(A\cup\{g\}\bigr)$.

\medskip
Relative to Theorem~\ref{thm:justifiable-choice}, we retain the same assumptions and add Strict One-Side
Monotonicity to recover beliefs and state-independent tastes from the
state-dependent representation. Although the auxiliary set of
state-dependent utilities is convex, utilities generated by
probability--utility pairs generally do not form a convex set: combining
utilities with different priors and utility functions need not preserve the
factorization $\nu(x,s)=\pi(s)u(x)$. As emphasized by Nau (2006) and
developed by Riella (2015) and Hara and Riella (2023), the state-independent
structure must instead be recovered from suitable extreme or exposed
generators of the convex set of state-dependent utilities.

We adapt this logic to choice correspondences. Strict One-Side Monotonicity plays
the same structural role here and, together with the axioms of Theorem~\ref{thm:justifiable-choice},
yields a justifiable representation in terms of
probability--utility pairs.

\begin{theorem}\label{thm:representation}
The following statements are equivalent:
\begin{enumerate}
    \item[(1)] The choice correspondence $c$ satisfies Axioms 1--7.
    \item[(2)] There exists a nonempty closed subset
    $\mathcal{M}\subseteq\Delta(S)\times C(X)$ of probability--utility pairs
    such that, for every menu $A\in\mathcal{A}$,
    \[
        c(A)=\bigcup_{(\pi,u)\in\mathcal{M}}
        \operatorname*{arg\,max}_{g\in A} V_{\pi,u}(g).
    \]
    Moreover, there exist $\hat p,\hat q\in\Delta(X)$ such that
    $\mathds{E}_{\hat p}(u)>\mathds{E}_{\hat q}(u)$
    for every $(\pi,u)\in\mathcal{M}$.
\end{enumerate}
\end{theorem}

The representation also shows how non-binary menus can separate
indecisiveness from indifference. Let $m=\lambda f+(1-\lambda)g$. Since each
$V_{\pi,u}$ is affine, $V_{\pi,u}(m)$ lies strictly between $V_{\pi,u}(f)$
and $V_{\pi,u}(g)$ unless $V_{\pi,u}(f)=V_{\pi,u}(g)$. A pair that strictly
ranks $f$ and $g$ therefore never makes $m$ optimal, and $m$ is chosen from
$\{f,g,m\}$ if and only if some pair is indifferent between $f$ and $g$.

Binary choice cannot tell the two cases apart: $f,g\in c(\{f,g\})$ whether
they are indifferent or ranked oppositely by different pairs. Adding $m$ to
the menu does. If $m$ is rejected, every pair has a strict ranking and the
two sides disagree, so the decision maker is indecisive between $f$ and $g$.
If $m$ is chosen, some pair is indifferent, though others may still disagree.
The test is thus conclusive for indecisiveness and only suggestive for
indifference.

\appendix
\section{Proofs} \label{sec:appendix}

\subsection{Proof of Proposition~\ref{prop:strict-representation}}

\begin{proof}[\textnormal{\textit{Necessity}}]
The necessity of the axioms is standard. Suppose that $\succ$ admits the
stated representation. Open Continuity
follows from the compactness of $\mathcal{V}$ and the continuity of
$U_\nu$. Independence follows from the affinity of each $U_\nu$, and
Constant Non-Triviality follows from the existence of $\hat p,\hat q$
such that $U_\nu(\hat p)>U_\nu(\hat q)$ for every $\nu\in\mathcal{V}$.
\end{proof}

\begin{proof}[\textnormal{\textit{Sufficiency}}]
Suppose that $\succ$ satisfies (S1)--(S3). By Constant Non-Triviality,
there exist $\hat p,\hat q\in\Delta(X)$ such that
$\hat p\succ\hat q$. Define the binary relation $\succsim^*$ on
$\mathcal{F}$ by
\[
    f\succsim^* g
    \iff
    \alpha\hat p+(1-\alpha)f
    \succ
    \alpha\hat q+(1-\alpha)g
    \quad\text{for every }\alpha\in(0,1).
\]

By the argument in the proof of Theorem 8 of Hara and Riella (2023),
$\succsim^*$ is a preorder satisfying Independence, Continuity, and
Directional Open-Continuity with respect to the pair
$(\hat p,\hat q)$. Moreover, $\succ\subseteq\succsim^*$.
We may therefore apply Theorem 1 of Hara and Riella (2023). In particular,
its proof yields a compact set $\mathcal{V}_0\subset C(X\times S)$ such
that
\[
    f\succsim^* g
    \iff
    U_\nu(f)\geq U_\nu(g)
    \quad\text{for every }\nu\in\mathcal{V}_0,
\]
and $U_\nu(\hat p)-U_\nu(\hat q)=1$ for every $\nu\in\mathcal{V}_0$.

We first show that $f\succ g$ implies $U_\nu(f)>U_\nu(g)$ for every
$\nu\in\mathcal{V}_0$.
Suppose that $f\succ g$. By Independence,
$f\succ \frac12 f+\frac12 g\succ g$.
Let $h=\frac12 f+\frac12 g$. By Open Continuity, there exist open
neighborhoods $O_f$ of $f$ and $O_g$ of $g$ such that $f'\succ h\succ g'$
for all $(f',g')\in O_f\times O_g$. By transitivity, $f'\succ g'$
for all $(f',g')\in O_f\times O_g$.

Since the mappings $\alpha\mapsto\alpha\hat q+(1-\alpha)f$
and $\alpha\mapsto\alpha\hat p+(1-\alpha)g$
are continuous and converge to $f$ and $g$, respectively, as
$\alpha\to 0$, there exists $\alpha\in(0,1)$ sufficiently small
such that $\alpha\hat q+(1-\alpha)f\in O_f$ and
$\alpha\hat p+(1-\alpha)g\in O_g$. Consequently,
$\alpha\hat q+(1-\alpha)f\succ\alpha\hat p+(1-\alpha)g$.
Since $\succ\subseteq\succsim^*$, the representation of
$\succsim^*$ implies that, for every $\nu\in\mathcal{V}_0$,
$\alpha U_\nu(\hat q)+(1-\alpha)U_\nu(f)
\geq\alpha U_\nu(\hat p)+(1-\alpha)U_\nu(g)$.
Rearranging, $U_\nu(f)-U_\nu(g)
\geq\frac{\alpha}{1-\alpha}
\left[U_\nu(\hat p)-U_\nu(\hat q)\right]$.
By the normalization obtained in the proof of Theorem 1 of Hara and Riella (2023), we have
$U_\nu(\hat p)-U_\nu(\hat q)=1$ for every $\nu\in\mathcal{V}_0$.
Therefore,
\[
    U_\nu(f)-U_\nu(g)
    \geq
    \frac{\alpha}{1-\alpha}>0
    \quad\text{for every }\nu\in\mathcal{V}_0.
\]
Hence, $f\succ g$ implies $U_\nu(f)>U_\nu(g)$ for every
$\nu\in\mathcal{V}_0$.

We now prove the converse implication. Suppose that
$f,g\in\mathcal{F}$ satisfy $U_\nu(f)>U_\nu(g)$ for every
$\nu\in\mathcal{V}_0$.
Since $\mathcal{V}_0$ is compact and the map
$\nu\mapsto U_\nu(f)-U_\nu(g)$ is continuous,
\[
    m:=\min_{\nu\in\mathcal{V}_0}
    \bigl[U_\nu(f)-U_\nu(g)\bigr]>0.
\]
Similarly, define $M:=\max_{\nu\in\mathcal{V}_0}
\bigl[U_\nu(\hat p)-U_\nu(\hat q)\bigr]$.
By the normalization obtained from Theorem 1, $M=1$.
Choose $\alpha\in(0,1)$ sufficiently small such that $(1-\alpha)m>\alpha M$.
Then, for every $\nu\in\mathcal{V}_0$, the affinity of $U_\nu$ implies
\[
\begin{aligned}
&U_\nu\bigl(\alpha\hat q+(1-\alpha)f\bigr)
-U_\nu\bigl(\alpha\hat p+(1-\alpha)g\bigr)\\
&\quad=(1-\alpha)
\bigl[U_\nu(f)-U_\nu(g)\bigr]
-\alpha
\bigl[U_\nu(\hat p)-U_\nu(\hat q)\bigr]\\
&\quad\geq(1-\alpha)m-\alpha M>0.
\end{aligned}
\]

Hence, by the representation of $\succsim^*$,
$\alpha\hat q+(1-\alpha)f\succsim^*\alpha\hat p+(1-\alpha)g$.
By the definition of $\succsim^*$, for every $\beta\in(0,1)$,
\[
\begin{aligned}
&\beta\hat p
+(1-\beta)\bigl[\alpha\hat q+(1-\alpha)f\bigr]\\
&\qquad\succ
\beta\hat q
+(1-\beta)\bigl[\alpha\hat p+(1-\alpha)g\bigr].
\end{aligned}
\]

In particular, let $\beta=\frac{\alpha}{1+\alpha}$.
Since $\beta=(1-\beta)\alpha$, simplifying the preceding comparison gives
\[
\begin{aligned}
&\frac{2\alpha}{1+\alpha}
\left(\frac{\hat p+\hat q}{2}\right)
+\frac{1-\alpha}{1+\alpha}f\\
&\qquad\succ
\frac{2\alpha}{1+\alpha}
\left(\frac{\hat p+\hat q}{2}\right)
+\frac{1-\alpha}{1+\alpha}g.
\end{aligned}
\]

Since $\frac{1-\alpha}{1+\alpha}\in(0,1)$, Independence implies $f\succ g$.
Therefore, $f\succ g\iff U_\nu(f)>U_\nu(g)$
for every $\nu\in\mathcal{V}_0$.

It remains to show that $\mathcal{V}$ can be chosen convex.
Let $\mathcal{V}_0$ be the compact representing set obtained
above and define $\mathcal{V}:=\overline{\operatorname{co}}(\mathcal{V}_0)$,
where the closure is taken in the supremum norm.
Since $C(X\times S)$ is a Banach space, the closed convex
hull of a compact set is compact. Moreover, for any
$f\succ g$, continuity and compactness imply
$m:=\min_{\nu\in\mathcal{V}_0}[U_\nu(f)-U_\nu(g)]>0$.
By linearity in $\nu$, this lower bound is preserved under
convex combinations and, by continuity, under closure.
Thus $U_\nu(f)>U_\nu(g)$ for every $\nu\in\mathcal{V}$.
The converse follows from $\mathcal{V}_0\subseteq\mathcal{V}$.
Hence, convexification preserves the representation.
\renewcommand{\qedsymbol}{$\blacksquare$}

\end{proof}

\subsection{Proof of Theorem~\ref{thm:justifiable-choice}}

\begin{proof}[Necessity]
All axioms except Continuity follow immediately from the representation
and the affinity of each $U_\nu$. We therefore verify only Continuity.

Let $A_n\to A$ in the Hausdorff metric and $f_n\to f$ in
$\mathcal F$, with $f_n\in c(A_n)$ for every $n$. For each $n$,
choose $\nu_n\in\mathcal W$ such that
$f_n\in\operatorname*{arg\,max}_{g\in A_n}U_{\nu_n}(g)$.
Compactness of $\mathcal W$ implies that, along a subsequence,
$\nu_n\to\nu\in\mathcal W$. Since $f_n\in A_n$ and $A_n\to A$,
we have $f\in A$. Fix $g\in A$. Hausdorff convergence yields
$g_n\in A_n$ with $g_n\to g$. Since $f_n$ maximizes $U_{\nu_n}$
on $A_n$, $U_{\nu_n}(f_n)\geq U_{\nu_n}(g_n)$ for every $n$.
By joint continuity of $(\nu,h)\mapsto U_\nu(h)$, passing to the
limit gives $U_\nu(f)\geq U_\nu(g)$. As $g\in A$ was arbitrary,
$f\in\operatorname*{arg\,max}_{g\in A}U_\nu(g)$, and hence
$f\in c(A)$.
\end{proof}

\begin{proof}[Sufficiency]
Suppose that Axioms 1--6 hold.
For any $f,g\in\mathcal{F}$, define the relation $\succ_c$ by
$f\succ_c g\Longleftrightarrow g\notin c(\{f,g\})$.

\begin{claim}\label{cl:binary-order}
The relation $\succ_c$ is a strict partial order that satisfies axioms
(S1)--(S3).
\end{claim}
\begin{proof}
Nonemptiness of $c(A)$ for every $A\in\mathcal{A}$, together with
Pseudo-WARP, implies that $\succ_c$ is a strict partial order.
Independence of $c(\cdot)$ implies that $\succ_c$ satisfies Independence.
Constant Non-Triviality of $c(\cdot)$ implies Constant Non-Triviality
of $\succ_c$: there exist $\hat p,\hat q\in\Delta(X)$ such that
$\hat q\notin c(\{\hat p,\hat q\})$, so $\hat p\succ_c\hat q$.
Similarly, Continuity of $c(\cdot)$ implies Open Continuity of $\succ_c$.
Therefore, $\succ_c$ satisfies all the properties required to apply
Proposition~\ref{prop:strict-representation}.
\end{proof}

\begin{claim}\label{cl:maximal-menu}
For every $A\in\mathcal{A}$ and $f\in c(A)$, the set
$\{B\in\mathcal{A}: A\subseteq B\text{ and }f\in c(B)\}$ has a maximal
element, $M$, with respect to set inclusion.
\end{claim}
\begin{proof}
Suppose $\mathcal{B}$ is a subset of
$\{B\in\mathcal{A}: A\subseteq B\text{ and }f\in c(B)\}$ linearly ordered
by $\supseteq$. Let $\bar{B}:=\operatorname{cl}(\bigcup\mathcal{B})$.
This implies that there exists a sequence $(B^m)$ of finite subsets of
$\bar{B}$ such that $B^m\to\bar{B}$. By the definition of $\bar{B}$, we
can find such a sequence with $B^m\subseteq\bigcup\mathcal{B}$ for every
$m\in\mathbb{N}$. Finally, since each $B^m$ is finite and $\mathcal{B}$
is linearly ordered by $\supseteq$, we can find, for each $m\in\mathbb{N}$,
a set $\bar{B}^m\in\mathcal{B}$ with $B^m\subseteq\bar{B}^m$. But then,
$\bar{B}^m\to\bar{B}$ and $f\in c(\bar{B}^m)$ for every
$m\in\mathbb{N}$. Continuity implies that $f\in c(\bar{B})$. That is,
$\bar{B}$ is an upper bound for $\mathcal{B}$. Now Zorn's lemma implies
that the collection
$\{B\in\mathcal{A}: A\subseteq B\text{ and }f\in c(B)\}$ has a maximal
element.
\end{proof}

\begin{claim}\label{cl:convex-menu}
Both $M$ and $\mathcal{F}\setminus M$ are convex.
\end{claim}
\begin{proof}
Since $f\in c(M)$, if $g,h\in M$ and $\alpha\in(0,1)$,
No-Hedging and maximality of $M$ immediately imply that
$\alpha g+(1-\alpha)h\in M$. Thus, $M$ is convex.

To prove that $\mathcal{F}\setminus M$ is convex, suppose that
$g\notin M$ and $h\notin M$, and let $\alpha\in(0,1)$.
By maximality of $M$, $f\notin c(M\cup\{g\})$ and
$f\notin c(M\cup\{h\})$. Therefore,
$f\notin c(M\cup\{g\})\cup c(M\cup\{h\})$.
Indirect Ambiguity Aversion then implies that
$f\notin c\bigl(M\cup\{\alpha g+(1-\alpha)h\}\bigr)$.
Since $f\in c(M)$, it follows that
$\alpha g+(1-\alpha)h\notin M$.
Hence, $\mathcal{F}\setminus M$ is convex.
\end{proof}

Fix $f\in\mathcal{F}$ and $A\in\mathcal{A}$ such that $f\in c(A)$.
Let $M$ be a maximal menu generated by this choice, as in
Claim~\ref{cl:maximal-menu}, so that $A\subseteq M$ and $f\in c(M)$.
Fix $\lambda\in(0,1)$ and a constant act $q\in\Delta(X)$ that returns a
full-support lottery in every state. Such a lottery exists because $X$
is compact metrizable.\footnote{A compact metric space has a countable
 dense sequence $(x_n)_{n\geq1}$. The lottery
 $q=\sum_{n=1}^{\infty}2^{-n}\delta_{x_n}$ assigns positive probability
 to every nonempty open subset of $X$ and therefore has full support.}
Define the operation $T_\lambda\colon\mathcal{F}\to\mathcal{F}_\lambda$
by $g\mapsto g_\lambda=\lambda g+(1-\lambda)q$, where
$\mathcal{F}_\lambda:=T_\lambda(\mathcal{F})$.
Observe that $T_\lambda$ is a homeomorphism onto $\mathcal{F}_\lambda$
and transforms general acts into acts that return full-support lotteries
in every state.
By Independence, $f\in c(M)$ implies $f_\lambda\in c(M_\lambda)$,
where $M_\lambda:=\lambda M+(1-\lambda)q$.
By Claim~\ref{cl:maximal-menu}, there exists a maximal menu
$M^*\in\mathcal{A}$ such that $M_\lambda\subseteq M^*$ and
$f_\lambda\in c(M^*)$. Additionally, we show that
$M_\lambda=M^*\cap\mathcal{F}_\lambda$.

The inclusion $M_\lambda\subseteq M^*\cap\mathcal{F}_\lambda$ follows
immediately from the construction of $M^*$ and the definition of
$\mathcal{F}_\lambda$. To establish the reverse inclusion, suppose that
$h\in M^*\cap\mathcal{F}_\lambda$. Then there exists $g\in\mathcal{F}$
such that $h=\lambda g+(1-\lambda)q=g_\lambda$.
Suppose, toward a contradiction, that $g_\lambda\notin M_\lambda$.
By the injectivity of $T_\lambda$, $g\notin M$.
Given that $f_\lambda\in c(M^*)$,
$f_\lambda\in M_\lambda\cup\{g_\lambda\}$, and
$c(M_\lambda\cup\{g_\lambda\})\subseteq M^*$, Pseudo-WARP implies
$f_\lambda\in c(M_\lambda\cup\{g_\lambda\})$.
Independence then gives $f\in c(M\cup\{g\})$.
By maximality of $M$, as in Claim~\ref{cl:maximal-menu}, this implies
$g\in M$, a contradiction. Therefore,
\begin{equation}\label{eq:transformed-maximal-menu}
 M_\lambda=M^*\cap\mathcal{F}_\lambda.
\end{equation}

We now apply the minimax theorem to obtain a state-dependent utility
$\nu^*$. By Claim~\ref{cl:convex-menu}, $M^*$ is convex, and it is
compact because $M^*\in\mathcal{A}$. By Claim~\ref{cl:binary-order}
and Proposition~\ref{prop:strict-representation}, let
$\mathcal{V}\subseteq C(X\times S)$ be a nonempty compact and convex
set representing $\succ_c$.
The function $U_\nu(f_\lambda)-U_\nu(g)$ is continuous and affine in
each variable, so Sion's Minimax Theorem applies to $M^*$ and
$\mathcal{V}$. By Pseudo-WARP, Claim~\ref{cl:binary-order}, and the
minimax theorem,
\[
\begin{aligned}
 f_\lambda\in c(M^*)
 &\implies \forall g\in M^*,\quad
 f_\lambda\in c(\{f_\lambda,g\})\\
 &\iff \forall g\in M^*,\quad g\not\succ_c f_\lambda\\
 &\iff \forall g\in M^*,\ \exists\nu\in\mathcal{V}:\quad
 U_\nu(f_\lambda)\geq U_\nu(g)\\
 &\iff \min_{g\in M^*}\max_{\nu\in\mathcal{V}}
 \bigl[U_\nu(f_\lambda)-U_\nu(g)\bigr]\geq0\\
 &\iff \max_{\nu\in\mathcal{V}}\min_{g\in M^*}
 \bigl[U_\nu(f_\lambda)-U_\nu(g)\bigr]\geq0\\
 &\iff \exists\nu^*\in\mathcal{V},\ \forall g\in M^*:\quad
 U_{\nu^*}(f_\lambda)\geq U_{\nu^*}(g).
\end{aligned}
\]
The extrema are attained by continuity and compactness.

From this result, we can also show that $M^*$ is the lower contour set
of $U_{\nu^*}$ with respect to $f_\lambda$. That is,
$M^*=\{g\in\mathcal{F}:U_{\nu^*}(g)\leq U^*\}$, where
$U^*:=U_{\nu^*}(f_\lambda)$. The inclusion
$M^*\subseteq\{g\in\mathcal{F}:U_{\nu^*}(g)\leq U^*\}$ follows
immediately from the application of the minimax theorem. For the other
direction, suppose, toward a contradiction, that $g\notin M^*$ but
$U_{\nu^*}(g)<U^*$. There exists an act $\tilde g\notin M^*$ and
a constant $C<\infty$ such that $\tilde g(s)\leq C f_\lambda(s)$ for
every $s\in S$ and $U_{\nu^*}(\tilde g)<U^*$.\footnote{Because each
$f_\lambda(s)$ has full support, the set of probabilities dominated by
a finite multiple of $f_\lambda(s)$ is weakly dense in $\Delta(X)$.
Indeed, finite-support probabilities are weakly dense, and each point
mass $\delta_x$ can be approximated by the normalized restrictions
$f_\lambda(s)(\,\cdot\cap B(x,\varepsilon))/
f_\lambda(s)(B(x,\varepsilon))$, which are dominated by a finite
multiple of $f_\lambda(s)$. Since $S$ is finite, the statewise constants
can be replaced by a common $C<\infty$. Hence, because
$\mathcal{F}\setminus M^*$ and
$\{h\in\mathcal{F}:U_{\nu^*}(h)<U^*\}$ are open, any $g$ in their intersection
can be approximated by some $\tilde g$ that remains in the intersection
and satisfies $\tilde g(s)\leq C f_\lambda(s)$ for every $s\in S$.}
The existence of $\tilde g$ also allows the construction of an act
$h=(1+t)f_\lambda-t\tilde g$ for sufficiently small
$t\in(0,1)$.\footnote{Choose $t\in(0,1)$ such that
$1+t-tC>0$. Then, for every $s\in S$,
$h(s)=(1+t)f_\lambda(s)-t\tilde g(s)$ is a nonnegative Borel measure
on $X$, since $h(s)\geq(1+t-tC)f_\lambda(s)\geq0$.
Moreover, $h(s)(X)=1$, so $h(s)\in\Delta(X)$.}
Thus $f_\lambda=\frac{t}{1+t}\tilde g+\frac{1}{1+t}h$ and
$U_{\nu^*}(h)=U^*+t[U^*-U_{\nu^*}(\tilde g)]>U^*$, so
$h\notin M^*$. Since $\tilde g\notin M^*$ as well, convexity of
$\mathcal{F}\setminus M^*$, by Claim~\ref{cl:convex-menu}, implies
$f_\lambda\notin M^*$, a contradiction.

This proves the first case. Now suppose that $g\notin M^*$ and
$U_{\nu^*}(g)=U^*$. Since $f_\lambda$ has full support, there exist
acts $g^+$ and $g^-$ such that
$U_{\nu^*}(g^-)<U^*<U_{\nu^*}(g^+)$. Indeed, by Constant
Non-Triviality and Proposition~\ref{prop:strict-representation},
$U_{\nu^*}$ is not constant on $\mathcal{F}$. Hence there exists
at least one state $s\in S$ for which $x\mapsto\nu^*(x,s)$ is
nonconstant. Since $X$ is compact and $\nu^*(\cdot,s)$ is continuous,
for every state $s$ there exist $x_s^-,x_s^+\in X$ attaining the minimum
and maximum of $\nu^*(\cdot,s)$, respectively.

Define $g^-,g^+\in\mathcal{F}$ statewise. Whenever
$\nu^*(\cdot,s)$ is nonconstant, let $g^-(s)=\delta_{x_s^-}$ and
$g^+(s)=\delta_{x_s^+}$. Whenever $\nu^*(\cdot,s)$ is constant, let
$g^-(s)=g^+(s)=f_\lambda(s)$. Because $f_\lambda(s)$ has full
support, in every state where $\nu^*(\cdot,s)$ is nonconstant we have
\[
 \nu^*(x_s^-,s)
 <\int_X\nu^*(x,s)\,df_\lambda(s)(x)
 <\nu^*(x_s^+,s).
\]
Indeed, the integral is a probability-weighted average of the values of
$\nu^*(\cdot,s)$, and full support assigns positive probability to
open sets on which $\nu^*(\cdot,s)$ is strictly above its minimum and
strictly below its maximum. In states where $\nu^*(\cdot,s)$ is
constant, the three quantities coincide. Since there is at least one
state in which $\nu^*(\cdot,s)$ is nonconstant, summing across states
yields
\[
 U_{\nu^*}(g^-)<U_{\nu^*}(f_\lambda)<U_{\nu^*}(g^+).
\]
Thus, we obtain $U_{\nu^*}(g^-)<U^*<U_{\nu^*}(g^+)$.

Now let $t\in(0,1)$ be sufficiently small and define the act
$h=t g^-+(1-t)g$ so that $h\notin M^*$. Since
$\mathcal{F}\setminus M^*$ is open and $h\to g$ as $t\downarrow0$,
such an act exists, and $U_{\nu^*}(h)<U^*$. But this is exactly the
contradiction established in the first case. Therefore,
\begin{equation}\label{eq:exact-transformed-contour}
 M^*=\{g\in\mathcal{F}:U_{\nu^*}(g)\leq U^*\}.
\end{equation}
By \eqref{eq:transformed-maximal-menu} and affinity, since $\lambda>0$,
\begin{equation}\label{eq:original-maximal-contour}
 M=\{g\in\mathcal{F}:U_{\nu^*}(g)\leq U_{\nu^*}(f)\}.
\end{equation}

Pseudo-WARP and Continuity also imply that if $g\in\mathcal{F}$
satisfies $U_{\nu^*}(g)=U^*$, then $g\in c(M^*)$.\footnote{To see
this, mix $g$ with $g^+$: for $t\in(0,1)$, let
$h=t g^++(1-t)g$. Then $U_{\nu^*}(h)>U^*$, so $h\notin M^*$.
Maximality implies $f_\lambda\notin c(M^*\cup\{h\})$.
If $h$ were also unchosen, then $c(M^*\cup\{h\})\subseteq M^*$;
Pseudo-WARP and $f_\lambda\in c(M^*)$ would imply
$f_\lambda\in c(M^*\cup\{h\})$, a contradiction. Hence
$h\in c(M^*\cup\{h\})$. As $t\downarrow0$, $h\to g$ and
$M^*\cup\{h\}\to M^*$, so Continuity gives $g\in c(M^*)$.}

Now fix $B\in\mathcal{A}$ and
$g\in\operatorname*{arg\,max}_{h\in B}U_{\nu^*}(h)$. We construct a
common mixture of $B$ whose $U_{\nu^*}$-maximizer lies on the boundary
level $U^*=U_{\nu^*}(f_\lambda)$. If $U_{\nu^*}(g)>U^*$, mix
$B$ with $g^-$ and choose $\alpha\in(0,1)$ so that
$U_{\nu^*}(\alpha g+(1-\alpha)g^-)=U^*$.\footnote{Take
$\alpha=(U^*-U_{\nu^*}(g^-))/(U_{\nu^*}(g)-U_{\nu^*}(g^-))$.}
If $U_{\nu^*}(g)<U^*$, mix $B$ with $g^+$ and choose
$\alpha\in(0,1)$ so that
$U_{\nu^*}(\alpha g+(1-\alpha)g^+)=U^*$.\footnote{Take
$\alpha=(U_{\nu^*}(g^+)-U^*)/(U_{\nu^*}(g^+)-U_{\nu^*}(g))$.}
If $U_{\nu^*}(g)=U^*$, take $\alpha=1$.

In every case, letting $j\in\{g^-,g^+\}$ denote the relevant anchor
(either may be used when $\alpha=1$), the optimality of $g$ on $B$ and
affinity of $U_{\nu^*}$ imply
$U_{\nu^*}(\alpha h+(1-\alpha)j)\leq U^*$ for every $h\in B$.
Hence $\alpha B+(1-\alpha)j\subseteq M^*$, while
$\alpha g+(1-\alpha)j$ lies on the boundary
$\{h\in\mathcal{F}:U_{\nu^*}(h)=U^*\}\subseteq c(M^*)$.
Pseudo-WARP and Independence then imply $g\in c(B)$. This proves that
\[
 \operatorname*{arg\,max}_{g\in B}U_{\nu^*}(g)\subseteq c(B)
 \quad\text{for every }B\in\mathcal{A}.
\]

Define
\[
 \mathcal W:=\left\{\nu\in\mathcal V:
 \operatorname*{arg\,max}_{g\in B}U_\nu(g)\subseteq c(B)
 \text{ for every }B\in\mathcal A\right\}.
\]
The preceding argument shows that $\mathcal W\neq\varnothing$.
Indeed, for every $A\in\mathcal A$ and every $f\in c(A)$, the
construction above produces some $\nu^*\in\mathcal V$ such that
$f\in\operatorname*{arg\,max}_{g\in A}U_{\nu^*}(g)$ and
$\operatorname*{arg\,max}_{g\in B}U_{\nu^*}(g)\subseteq c(B)$ for
every $B\in\mathcal A$. Hence $\nu^*\in\mathcal W$. It follows
immediately that, for every $A\in\mathcal A$,
\[
 c(A)=\bigcup_{\nu\in\mathcal W}
 \operatorname*{arg\,max}_{g\in A}U_\nu(g).
\]

It remains to show that $\mathcal W$ is compact. Since
$\mathcal W\subseteq\mathcal V$ and $\mathcal V$ is compact, it
suffices to prove that $\mathcal W$ is closed. Let
$\nu_n\in\mathcal W$ and $\nu_n\to\nu\in\mathcal V$ in the
supremum norm. Fix $B\in\mathcal A$ and
$g\in\operatorname*{arg\,max}_{h\in B}U_\nu(h)$.

By Proposition~\ref{prop:strict-representation}, there are constant
acts $\hat p,\hat q$ such that $U_\mu(\hat p)>U_\mu(\hat q)$ for
every $\mu\in\mathcal V$. Choose $\beta\in(0,1)$ small enough that,
for $B_\beta:=\beta B+(1-\beta)\hat q$ and
$g_\beta:=\beta g+(1-\beta)\hat q$, we have
$U_\nu(g_\beta)<U_\nu(\hat p)$. Affinity implies
$g_\beta\in\operatorname*{arg\,max}_{h\in B_\beta}U_\nu(h)$.

For $t\in(0,1)$, let $g_\beta^t:=(1-t)g_\beta+t\hat p$.
Then, for every $h\in B_\beta$,
\[
 U_\nu(g_\beta^t)-U_\nu(h)
 \geq t\bigl[U_\nu(\hat p)-U_\nu(g_\beta)\bigr]>0.
\]
Since $\nu_n\to\nu$ uniformly, for each fixed $t$ the difference
remains strictly positive under $U_{\nu_n}$, uniformly over
$h\in B_\beta$, for all sufficiently large $n$. Hence
$g_\beta^t\in c(B_\beta\cup\{g_\beta^t\})$, because
$\nu_n\in\mathcal W$. Letting $t\downarrow0$, Continuity gives
$g_\beta\in c(B_\beta)$, and Independence then yields $g\in c(B)$.

Thus every $U_\nu$-maximizer is chosen, so $\nu\in\mathcal W$.
Therefore $\mathcal W$ is closed in $\mathcal V$, hence compact.
Since $\mathcal W\subseteq\mathcal V$,
$U_\nu(\hat p)>U_\nu(\hat q)$ for every $\nu\in\mathcal W$.
\renewcommand{\qedsymbol}{$\blacksquare$}
\end{proof}

\singlespacing
\setlength\bibsep{0.8em}
\begin{thebibliography}{99}

\bibitem[Aizerman and Malishevski(1981)]{aizerman1981}
Aizerman, M., and A. Malishevski (1981).
General theory of best variants choice: Some aspects.
\textit{IEEE Transactions on Automatic Control}, 26(5), 1030--1040.

\bibitem[Eliaz and Ok(2006)]{eliaz2006}
Eliaz, K., and E. A. Ok (2006).
Indifference or indecisiveness? Choice-theoretic foundations of incomplete preferences.
\textit{Games and Economic Behavior}, 56(1), 61--86.

\bibitem[Galaabaatar and Karni(2012)]{galaabaatar2012}
Galaabaatar, T., and E. Karni (2012).
Expected multi-utility representations.
\textit{Mathematical Social Sciences}, 64(3), 242--246.

\bibitem[Galaabaatar and Karni(2013)]{galaabaatar2013}
Galaabaatar, T., and E. Karni (2013).
Subjective expected utility with incomplete preferences.
\textit{Econometrica}, 81(1), 255--284.

\bibitem[Hara and Riella(2023)]{hara2023}
Hara, K., and G. Riella (2023).
Multiple tastes and beliefs with an infinite prize space.
\textit{Economic Theory}, 76(2), 417--444.

\bibitem[Heller(2012)]{heller2012}
Heller, Y. (2012).
Justifiable choice.
\textit{Games and Economic Behavior}, 76(2), 375--390.

\bibitem[Nau(2006)]{nau2006}
Nau, R. (2006).
The shape of incomplete preferences.
\textit{The Annals of Statistics}, 34(5), 2430--2448.

\bibitem[Ok et~al.(2012)Ok, Ortoleva, and Riella]{ok2012}
Ok, E. A., P. Ortoleva, and G. Riella (2012).
Incomplete preferences under uncertainty: Indecisiveness in beliefs versus tastes.
\textit{Econometrica}, 80(4), 1791--1808.

\bibitem[Riella(2015)]{riella2015}
Riella, G. (2015).
On the representation of incomplete preferences under uncertainty with indecisiveness in tastes and beliefs.
\textit{Economic Theory}, 58(3), 571--600.

\bibitem[Safonov(2023)]{safonov}
Safonov, E. (2023.).
Framing of ambiguity.
Available at SSRN: \url{https://ssrn.com/abstract=5260555}.

\end{thebibliography}

\end{document}
